\bookchapter{Capítulo 67: Refinamento e publicação — testes, depuração e publicação na App Store}

\chapterintro{O App ficou pronto, e agora? Este capítulo responde às três coisas que levam de “está rodando” a “está pronto para entregar”: como depurar e achar bugs com eficiência, como escrever testes para evitar regressões e como levar o App à App Store para o mundo inteiro baixar. Publicar não é só cerimônia — é a linha que separa uma obra de um produto.}

\booksection{Os três instrumentos de depuração: print, breakpoints e preview}

\begin{center}
\begin{tabularx}{\linewidth}{|X|X|X|}
\hline
\textbf{Ferramenta} & \textbf{Para que serve} & \textbf{Como usar} \tabularnewline
\hline
print + console & Ver logs de execução & No código, print("valor=\textbackslash{}(x)") e veja a saída no console na parte inferior do Xcode \tabularnewline
Depuração com breakpoints & Fazer o programa parar em uma linha e ver os valores das variáveis & Clique ao lado do número da linha para criar a marca azul; a execução pausa ao chegar ali e você avança passo a passo \tabularnewline
Preview do SwiftUI + depurador & Ver a interface e a hierarquia de views em tempo real & Canvas do \#Preview + inspetor View Hierarchy em tempo de execução \tabularnewline
\hline
\end{tabularx}
\end{center}

\begin{center}\small\color{gray}Os três instrumentos de depuração\end{center}

\begin{tcbinfo}{Dica: como ler um crash?}
Quando o Swift trava, o console imprime “Fatal error” seguido do motivo, e o Xcode para na linha do erro e a destaca. O histórico de travamentos em aparelhos reais pode ser visto em Xcode Window Organizer Crashes (requer conta paga). Priorize pesquisar “primeira frase do motivo do erro + swiftui” — só uma pergunta precisa gera uma resposta precisa.

\end{tcbinfo}

\booksection{Testes: um seguro para o código}

O teste manual não é confiável — você mexe em uma função e outra quebra às escondidas. Testes são o “remendo do arrependimento” e o “colete à prova de balas” ao mesmo tempo. No iOS usa-se o framework XCTest, em dois tipos:

\begin{center}
\begin{tabularx}{\linewidth}{|X|X|X|X|}
\hline
\textbf{Nível} & \textbf{O que testa} & \textbf{Target} & \textbf{Velocidade} \tabularnewline
\hline
Teste unitário & Lógica pura (cálculos, validações) & MyFirstAppTests & Segundos \tabularnewline
Teste de UI & Comportamento real da interface & MyFirstAppUITests & Minutos \tabularnewline
Teste manual & Sensação e experiência de pessoas reais & —— & Sempre necessário \tabularnewline
\hline
\end{tabularx}
\end{center}

\begin{center}\small\color{gray}Níveis de teste no iOS (o projeto-modelo já vem com dois Test Targets)\end{center}

\begin{lstlisting}[language=Swift]
import XCTest

// Extraia lógica como "string de valor para número" numa função pura
func parseAmount(_ input: String) -> Double? {
    guard let v = Double(input.trimmingCharacters(in: .whitespaces)),
          v > 0 else { return nil }
    return v
}

final class AmountTests: XCTestCase {
    func testEmptyInputReturnsNil() {
        XCTAssertNil(parseAmount(""))
    }
    func testNegativeReturnsNil() {
        XCTAssertNil(parseAmount("-5"))
    }
    func testNormalValueParses() {
        XCTAssertEqual(parseAmount("12.5"), 12.5)
    }
}
\end{lstlisting}

Rodar os testes: CmdU. Depois que está tudo verde, mexer no código fica muito mais tranquilo. A validação de valores e os cálculos de resumo do livro-caixa merecem testes unitários protegendo tudo.

\booksection{Assinatura e empacotamento: carimbando o App}

Todo App iOS precisa ser assinado para poder ser instalado — é a prova de que “este pacote veio de você e não foi adulterado”. A hierarquia é a seguinte:

\begin{itemize}
  \item Fase de aprendizado: Apple ID gratuito + assinatura automática do Xcode (Personal Team); roda em aparelho real, com renovação a cada 7 dias;
  \item Fase de distribuição: a conta de desenvolvedor paga (99 dólares/ano) emite os certificados e os perfis de provisionamento, necessários para publicar e instalar a longo prazo.
\end{itemize}

Empacotar: no menu do Xcode, Product Archive; após o build, no Organizer use Distribute App escolha App Store Connect faça o upload. O archive precisa ter o Generic iOS Device como destino (não o simulador) — um passo que os iniciantes costumam esquecer.

\textbf{Alerta de armadilha: perdeu o certificado = atualizações paradas}\par

Se o certificado de assinatura e a chave privada se perderem, você não consegue mais publicar atualizações do App antigo — os usuários só conseguiriam instalar o aplicativo como se fosse “novo”. Faça backup do certificado com cuidado (exporte o .p12 do keychain) e guarde bem a senha. Para ajudar, a Apple oferece a “assinatura automática” e o gerenciamento de certificados na nuvem (marque Automatically manage signing no Xcode), o que elimina a maioria das armadilhas manuais.

\booksection{O processo completo de publicação na App Store}

Registrar conta de desenvolvedor

Preparar materiais da loja

Enviar o build

Preencher informações de privacidade e revisão

Revisão da Apple (1-7 dias)

Publicado

\begin{center}\small\color{gray}Figura 5: os seis passos para publicar na App Store\end{center}

\begin{itemize}
  \item Cadastro: developer.apple.com, conta individual por 99 dólares/ano;
  \item Materiais: ícone (1024×1024, sem cantos arredondados — a Apple recorta sozinha), capturas de tela (nos vários tamanhos de iPhone), descrição e link da política de privacidade;
  \item Questionário de privacidade: declare com honestidade quais dados são coletados — um livro-caixa que não se conecta à internet pode marcar "não coleta nenhum dado";
  \item Revisão: a primeira submissão costuma levar de 1 a 7 dias de espera; se for rejeitado, corrija item a item conforme as "notas da revisão" — a Apple costuma escrever os motivos da rejeição de forma bem específica, basta seguir;
  \item TestFlight: antes de publicar, envie a versão de teste para amigos (até 10.000 pessoas), colete um ciclo de feedback e só então submeta de verdade — postura bem profissional.
\end{itemize}

\booksection{Como criar os materiais da loja: a primeira impressão vale downloads}

A página da loja é a vitrine que o usuário vê ao “passar em frente à loja”. Três materiais definem a primeira impressão:

\begin{itemize}
  \item Ícone: 1024×1024; não vale improvisar com uma captura de tela. Se você não tem habilidade de design, use uma ferramenta de IA para gerar uma forma simples; ele só passa no teste se continuar reconhecível de relance mesmo reduzido na tela inicial;
  \item Capturas de tela: pelo menos 2 conjuntos (tamanhos principais do iPhone), mostrando os cenários de uso mais centrais; cada uma com uma frase em linguagem humana ("Anote um gasto em um segundo" é dez vezes melhor do que "Solução eficiente de gestão financeira");
  \item Descrição: a primeira frase deve deixar claro "o que é + para quem é"; nada de empilhar "o mais profissional, o mais poderoso" — sinceridade funciona melhor do que exagero.
\end{itemize}

\begin{tcbpractice}{Mãos à obra: "entregue" o livro-caixa uma vez}
\booksection{Relatório de crashes e mini checklist de performance}

Antes de publicar o App, vale dedicar meia hora a cada uma destas duas tarefas:

\booksubsection{Monitoramento de crashes}

Não travar nos seus testes não significa que não vai travar no celular dos usuários. A Apple oferece relatórios de crash gratuitos no Xcode Organizer (contas pagas veem logs de crash de aparelhos reais); uma solução mais prática é integrar um SDK de terceiros (como o Firebase Crashlytics). Publicar sem monitoramento de crashes é dirigir de olhos fechados.

\booksubsection{Quatro checagens de performance}

\begin{enumerate}
  \item Listas com LazyVStack / List: listas longas não devem ser montadas com um VStack comum — carregamento preguiçoso é a primeira lição de performance;
  \item Não force as imagens originais: use AsyncImage (imagens remotas) e o cache do sistema; uma imagem original de 5 MB é capaz de travar aparelhos antigos;
  \item Não faça trabalho pesado na thread principal: rede e arquivos grandes devem passar por async/await;
  \item Meça com o Instruments: o analisador de performance que vem no Xcode (Time Profiler, checagem de vazamento de memória) — primeiro meça, depois otimize; não confie no feeling.
\end{enumerate}

(1) Adicione 3 testes unitários para lógicas como parseAmount (rodando verdes com CmdU); (2) use um breakpoint para caçar com as próprias mãos um crash que você mesmo criou e conserte-o; (3) gere um pacote com Product Archive e convide amigos para testar pelo TestFlight; (4) se sobrar energia, registre a conta de desenvolvedor e submeta o livro-caixa de verdade para publicação. Os três primeiros passos dão para fazer esta noite.

\end{tcbpractice}

\begin{tcbsummary}{Resumo do capítulo}
\begin{itemize}
  \item Três instrumentos de depuração: logs com print, breakpoints, preview do SwiftUI e inspetor de views
  \item XCTest: teste unitário (CmdU) + teste de UI + teste manual, já pré-configurado no modelo
  \item Assinatura: conta gratuita com assinatura automática (renovação em 7 dias); só a conta paga permite distribuir
  \item Empacotamento: Generic iOS Device + Archive + Distribute; faça backup do certificado sem falta
  \item Seis passos para publicar: cadastro materiais upload questionário de privacidade revisão publicação; teste antes com o TestFlight
\end{itemize}

\end{tcbsummary}
